\begin{theorembox}{Theorem (Group Cancellation)}
\begin{theorem}[Group Cancellation]
\label{thm:group-cancellation}
In a group, cancellation holds on both sides:
\[
a\ast c=b\ast c\Longrightarrow a=b
\qquad\text{and}\qquad
c\ast a=c\ast b\Longrightarrow a=b.
\]
\end{theorem}
\end{theorembox}

\begin{dependencies}
\begin{itemize}
  \item \hyperref[def:group]{Group}
  \item \hyperref[thm:uniqueness-of-inverse]{Uniqueness of Inverse}
\end{itemize}
\end{dependencies}

\begin{theorembox}{Theorem (Group Uniqueness of Inverse)}
\begin{theorem}[Group Uniqueness of Inverse]
\label{thm:group-uniqueness-of-inverse}
In a group, each element has a unique inverse.
\end{theorem}
\end{theorembox}

\begin{dependencies}
\begin{itemize}
  \item \hyperref[def:group]{Group}
  \item \hyperref[thm:uniqueness-of-inverse]{Uniqueness of Inverse}
\end{itemize}
\end{dependencies}

\begin{theorembox}{Theorem (Field Zero Product)}
\begin{theorem}[Field Zero Product]
\label{thm:field-zero-product}
In a field,
\[
a\cdot b=0
\quad\Longleftrightarrow\quad
a=0\ \text{or}\ b=0.
\]
\end{theorem}
\end{theorembox}

\begin{dependencies}
\begin{itemize}
  \item \hyperref[def:abstract-field]{Field}
  \item \hyperref[thm:no-zero-divisors]{No Zero Divisors}
\end{itemize}
\end{dependencies}

\begin{theorembox}{Theorem (Ring Double Negation)}
\begin{theorem}[Ring Double Negation]
\label{thm:ring-double-negation}
In a ring,
\[
-(-a)=a.
\]
\end{theorem}
\end{theorembox}

\begin{dependencies}
\begin{itemize}
  \item \hyperref[def:ring]{Ring}
  \item \hyperref[thm:uniqueness-of-inverse]{Uniqueness of Inverse}
\end{itemize}
\end{dependencies}

\begin{theorembox}{Theorem (Ring Negation of Product)}
\begin{theorem}[Ring Negation of Product]
\label{thm:ring-negation-of-product}
In a ring,
\[
(-a)\cdot b=-(a\cdot b)
\qquad\text{and}\qquad
a\cdot(-b)=-(a\cdot b).
\]
\end{theorem}
\end{theorembox}

\begin{dependencies}
\begin{itemize}
  \item \hyperref[prop:sign-rule]{Sign Rule}
\end{itemize}
\end{dependencies}

\begin{theorembox}{Theorem (Ring Product of Negatives)}
\begin{theorem}[Ring Product of Negatives]
\label{thm:ring-product-of-negatives}
In a ring,
\[
(-a)\cdot(-b)=a\cdot b.
\]
\end{theorem}
\end{theorembox}

\begin{dependencies}
\begin{itemize}
  \item \hyperref[prop:product-of-negatives]{Product of Negatives}
\end{itemize}
\end{dependencies}

\begin{theorembox}{Theorem (Field Multiplicative Cancellation)}
\begin{theorem}[Field Multiplicative Cancellation]
\label{thm:field-multiplicative-cancellation}
In a field,
\[
c\ne0\land a\cdot c=b\cdot c
\Longrightarrow
a=b.
\]
\end{theorem}
\end{theorembox}

\begin{dependencies}
\begin{itemize}
  \item \hyperref[def:abstract-field]{Field}
  \item \hyperref[thm:group-cancellation]{Group Cancellation}
\end{itemize}
\end{dependencies}

\begin{theorembox}{Theorem (Field Uniqueness of Inverse)}
\begin{theorem}[Field Uniqueness of Inverse]
\label{thm:field-uniqueness-of-inverse}
In a field, each nonzero element has a unique multiplicative inverse.
\end{theorem}
\end{theorembox}

\begin{dependencies}
\begin{itemize}
  \item \hyperref[def:abstract-field]{Field}
  \item \hyperref[thm:group-uniqueness-of-inverse]{Group Uniqueness of Inverse}
\end{itemize}
\end{dependencies}

\begin{theorembox}{Theorem (Field Inverse of One)}
\begin{theorem}[Field Inverse of One]
\label{thm:field-inverse-of-one}
In a field,
\[
1^{-1}=1.
\]
\end{theorem}
\end{theorembox}

\begin{dependencies}
\begin{itemize}
  \item \hyperref[thm:field-uniqueness-of-inverse]{Field Uniqueness of Inverse}
\end{itemize}
\end{dependencies}

\begin{theorembox}{Theorem (Field Inverse of Product)}
\begin{theorem}[Field Inverse of Product]
\label{thm:field-inverse-of-product}
In a field, if $a\ne0$ and $b\ne0$, then
\[
(a\cdot b)^{-1}=a^{-1}\cdot b^{-1}.
\]
\end{theorem}
\end{theorembox}

\begin{dependencies}
\begin{itemize}
  \item \hyperref[thm:field-uniqueness-of-inverse]{Field Uniqueness of Inverse}
  \item \hyperref[thm:field-zero-product]{Field Zero Product}
\end{itemize}
\end{dependencies}

\begin{theorembox}{Theorem (Field Inverse Involution)}
\begin{theorem}[Field Inverse Involution]
\label{thm:field-inverse-involution}
In a field, if $a\ne0$, then
\[
(a^{-1})^{-1}=a.
\]
\end{theorem}
\end{theorembox}

\begin{dependencies}
\begin{itemize}
  \item \hyperref[thm:field-uniqueness-of-inverse]{Field Uniqueness of Inverse}
\end{itemize}
\end{dependencies}

\begin{definitionbox}{Definition (Field Integer Power)}
\begin{definition}[Field Integer Power]
\label{def:field-integer-power}
Let $F$ be a field. For $a\in F$, powers $a^n$ with $n\in\mathbb{W}$ are
defined recursively. If $a\ne0$ and $n\in\mathbb{Z}$ is negative, define
\[
a^n:=(a^{-1})^{-n}.
\]
\end{definition}
\end{definitionbox}

\begin{dependencies}
\begin{itemize}
  \item \hyperref[def:abstract-field]{Field}
  \item \hyperref[thm:field-uniqueness-of-inverse]{Field Uniqueness of Inverse}
\end{itemize}
\end{dependencies}

\begin{theorembox}{Theorem (Field Exponent Addition)}
\begin{theorem}[Field Exponent Addition]
\label{thm:field-exponent-addition}
In a field, if $a\ne0$ and $m,n\in\mathbb{Z}$, then
\[
a^{m+n}=a^m\cdot a^n.
\]
\end{theorem}
\end{theorembox}

\begin{dependencies}
\begin{itemize}
  \item \hyperref[def:field-integer-power]{Field Integer Power}
\end{itemize}
\end{dependencies}

\begin{theorembox}{Theorem (Field Power of Power)}
\begin{theorem}[Field Power of Power]
\label{thm:field-power-of-power}
In a field, if $a\ne0$ and $m,n\in\mathbb{Z}$, then
\[
(a^m)^n=a^{m\cdot n}.
\]
\end{theorem}
\end{theorembox}

\begin{dependencies}
\begin{itemize}
  \item \hyperref[def:field-integer-power]{Field Integer Power}
\end{itemize}
\end{dependencies}

\begin{theorembox}{Theorem (Field Power of Product)}
\begin{theorem}[Field Power of Product]
\label{thm:field-power-of-product}
In a field, if $a\ne0$, $b\ne0$, and $n\in\mathbb{Z}$, then
\[
(a\cdot b)^n=a^n\cdot b^n.
\]
\end{theorem}
\end{theorembox}

\begin{dependencies}
\begin{itemize}
  \item \hyperref[def:field-integer-power]{Field Integer Power}
  \item \hyperref[thm:field-inverse-of-product]{Field Inverse of Product}
\end{itemize}
\end{dependencies}

\begin{corollarybox}{Corollary (Generic Field Laws)}
\begin{corollary}[Generic Field Laws]
\label{cor:generic-field-laws}
Every field satisfies the zero-product law, cancellation laws, uniqueness of
inverses, inverse laws, ring sign laws, and integer-exponent laws listed in
this topic.
\end{corollary}
\end{corollarybox}

\begin{dependencies}
\begin{itemize}
  \item \hyperref[thm:field-zero-product]{Field Zero Product}
  \item \hyperref[thm:field-multiplicative-cancellation]{Field Multiplicative Cancellation}
  \item \hyperref[thm:field-inverse-of-product]{Field Inverse of Product}
  \item \hyperref[thm:field-exponent-addition]{Field Exponent Addition}
\end{itemize}
\end{dependencies}

\begin{theorembox}{Theorem (Ordered Field Squares Are Nonnegative)}
\begin{theorem}[Ordered Field Squares Are Nonnegative]
\label{thm:ordered-field-squares-nonnegative}
In an ordered field,
\[
0\le a^2.
\]
\end{theorem}
\end{theorembox}

\begin{dependencies}
\begin{itemize}
  \item \hyperref[def:abstract-ordered-field]{Ordered Field}
\end{itemize}
\end{dependencies}

\begin{theorembox}{Theorem (Ordered Field One Is Positive)}
\begin{theorem}[Ordered Field One Is Positive]
\label{thm:ordered-field-one-positive}
In an ordered field,
\[
0<1.
\]
\end{theorem}
\end{theorembox}

\begin{dependencies}
\begin{itemize}
  \item \hyperref[def:abstract-ordered-field]{Ordered Field}
  \item \hyperref[thm:ordered-field-squares-nonnegative]{Ordered Field Squares Are Nonnegative}
\end{itemize}
\end{dependencies}

\begin{theorembox}{Theorem (Ordered Field Positive Product)}
\begin{theorem}[Ordered Field Positive Product]
\label{thm:ordered-field-positive-product}
In an ordered field,
\[
0<a\land0<b
\Longrightarrow
0<a\cdot b.
\]
\end{theorem}
\end{theorembox}

\begin{dependencies}
\begin{itemize}
  \item \hyperref[def:abstract-ordered-field]{Ordered Field}
\end{itemize}
\end{dependencies}

\begin{theorembox}{Theorem (Ordered Field Positive Inverses)}
\begin{theorem}[Ordered Field Positive Inverses]
\label{thm:ordered-field-positive-inverses}
In an ordered field,
\[
0<a
\Longrightarrow
0<a^{-1}.
\]
\end{theorem}
\end{theorembox}

\begin{dependencies}
\begin{itemize}
  \item \hyperref[thm:ordered-field-positive-product]{Ordered Field Positive Product}
  \item \hyperref[thm:field-inverse-of-one]{Field Inverse of One}
\end{itemize}
\end{dependencies}

\begin{theorembox}{Theorem (Ordered Field Positive Multiplication Preserves Order)}
\begin{theorem}[Ordered Field Positive Multiplication Preserves Order]
\label{thm:ordered-field-positive-multiplication-preserves-order}
In an ordered field,
\[
0<c\land a<b
\Longrightarrow
c\cdot a<c\cdot b.
\]
The corresponding non-strict implication also holds.
\end{theorem}
\end{theorembox}

\begin{dependencies}
\begin{itemize}
  \item \hyperref[def:abstract-ordered-field]{Ordered Field}
  \item \hyperref[thm:ordered-field-positive-product]{Ordered Field Positive Product}
\end{itemize}
\end{dependencies}

\begin{theorembox}{Theorem (Ordered Field Negative Multiplication Reverses Order)}
\begin{theorem}[Ordered Field Negative Multiplication Reverses Order]
\label{thm:ordered-field-negative-multiplication-reverses-order}
In an ordered field,
\[
c<0\land a<b
\Longrightarrow
c\cdot b<c\cdot a.
\]
The corresponding non-strict implication also holds.
\end{theorem}
\end{theorembox}

\begin{dependencies}
\begin{itemize}
  \item \hyperref[thm:ordered-field-positive-multiplication-preserves-order]{Ordered Field Positive Multiplication Preserves Order}
  \item \hyperref[thm:ring-negation-of-product]{Ring Negation of Product}
\end{itemize}
\end{dependencies}

\begin{theorembox}{Theorem (Ordered Field Fraction Comparison)}
\begin{theorem}[Ordered Field Fraction Comparison]
\label{thm:ordered-field-fraction-comparison}
In an ordered field, if $b\cdot d>0$, then
\[
\frac{a}{b}<\frac{c}{d}
\quad\Longleftrightarrow\quad
a\cdot d<c\cdot b.
\]
\end{theorem}
\end{theorembox}

\begin{dependencies}
\begin{itemize}
  \item \hyperref[thm:ordered-field-positive-multiplication-preserves-order]{Ordered Field Positive Multiplication Preserves Order}
  \item \hyperref[thm:field-multiplicative-cancellation]{Field Multiplicative Cancellation}
\end{itemize}
\end{dependencies}

\begin{corollarybox}{Corollary (Generic Ordered-Field Laws)}
\begin{corollary}[Generic Ordered-Field Laws]
\label{cor:generic-ordered-field-laws}
Every ordered field satisfies the standard positivity, order-compatibility,
inverse-order, sign, and comparison laws listed in this topic.
\end{corollary}
\end{corollarybox}

\begin{dependencies}
\begin{itemize}
  \item \hyperref[thm:ordered-field-one-positive]{Ordered Field One Is Positive}
  \item \hyperref[thm:ordered-field-positive-product]{Ordered Field Positive Product}
  \item \hyperref[thm:ordered-field-positive-inverses]{Ordered Field Positive Inverses}
  \item \hyperref[thm:ordered-field-positive-multiplication-preserves-order]{Ordered Field Positive Multiplication Preserves Order}
\end{itemize}
\end{dependencies}
